\documentclass[11pt]{article}
\usepackage[utf8]{inputenc}
\usepackage[T1]{fontenc}
\usepackage[spanish]{babel}
\usepackage{hyperref}
\title{TerraPulse AI: interfaz en español e inglés}
\author{TerraPulse AI}
\date{\today}
\begin{document}
\maketitle
\section{Objetivo}
Ofrecer la interfaz web en español e inglés con terminología de observación
terrestre, radar y simulación. La localización no modifica mediciones,
geometrías, identificadores científicos, algoritmos ni solicitudes de la API.
No traduce el código fuente, las herramientas de terminal ni los datos originales.
\section{Uso}
Seleccione Español o English en la barra superior, junto a Archivo,
Escenas GCOV y Procesamiento SAR. El español es el idioma inicial cuando no
existe una preferencia válida. El cambio actualiza menús, paneles, mapa,
gráficas, visor científico, búsquedas, catálogo, ayudas, avisos y errores conocidos.
La preferencia utiliza exclusivamente la clave
\texttt{terrapulse:language:v1} del almacenamiento local del navegador;
no reemplaza los mapas guardados. El idioma del documento HTML también se actualiza.
Si el navegador bloquea el almacenamiento, la selección sigue funcionando
durante la sesión, pero no se garantiza su conservación después de recargar.
\section{Implementación}
\texttt{web/src/i18n/core.ts} administra idioma, suscripciones, interpolación
y traducciones contextuales explícitas. Los catálogos separados agrupan textos
de navegación, mapa, gráficas y simulaciones. Los mensajes conocidos admiten
traducción de ida y vuelta, incluso si fueron almacenados mientras se mostraba inglés.
Los parámetros interpolados conservan números, identificadores y nombres originales.
\texttt{useTranslation.ts} utiliza \texttt{useSyncExternalStore} para actualizar
componentes sin volver a montar el mapa. Las opciones de Chart.js reciben la
configuración regional seleccionada. Los controles MapLibre actualizan su texto
sin cambiar la vista ni iniciar otra consulta por un cambio de idioma.
El cliente de la API localiza los mensajes conocidos de \texttt{detail} y de
validación de campos. Un cuerpo vacío, un mensaje vacío o JSON inválido conservan
el mensaje de respaldo con el estado HTTP. Las respuestas exitosas y los cuerpos
de solicitud permanecen intactos. No se agregan dependencias.
\subsection{Terminología contextual}
\begin{tabular}{ll}
Español & Inglés \\
\hline
Retrodispersión & Backscatter \\
Adquisición & Acquisition \\
Región de interés & Region of interest (ROI) \\
Trayectoria y marco & Track and frame \\
Madurez del producto & Data maturity \\
Calado & Water depth \\
Caudal & Discharge \\
Escorrentía & Runoff \\
Irradiancia incidente & Incident irradiance \\
\end{tabular}
Las unidades, polarizaciones, GCOV, NISAR, CRID, EPSG y fórmulas mantienen
su significado. Una observación no se convierte en una inferencia o predicción;
se conservan las advertencias de resultados sintéticos, exploratorios y no validados.
\section{Limitaciones}
Los nombres propios y el contenido introducido por usuarios no se traducen.
Las imágenes PNG científicas originales conservan el texto incrustado; la
interfaz del visor, sus títulos y descripciones sí cambian de idioma.
Los mensajes externos futuros que no existan en los catálogos se muestran tal
como llegan, en lugar de inventar una interpretación.
La traducción no establece exactitud hidráulica ni validación científica.
No se realizó una inspección visual o interacción en navegador: no hay un
navegador disponible en este entorno. No se ejecutan servicios Earthdata reales.
\section{Verificación}
Desde \texttt{web/}, ejecute \texttt{node --test tests/*.test.mjs}.
Último resultado observado: 83 pruebas aprobadas, cero fallidas.
Se comprueban renderizados públicos en ambos idiomas, preservación de
GeoJSON con huecos y múltiples partes, solicitudes, respuestas, identificadores,
unidades, datos de gráficas, persistencia separada y mensajes reversibles.
La prueba de las ecuaciones plegadas combina comprobación estructural con
traducción pública; no equivale a abrir el control en un navegador.
\texttt{npm run build} terminó con código cero también después de los
últimos ajustes de mensajes. Incluye compilación TypeScript y producción Vite.
Vite advirtió sobre \texttt{\_\_dirname} en una futura configuración y sobre
el tamaño de un fragmento JavaScript; no se modificó la configuración para ocultarlas.
\texttt{npm run lint} terminó con código uno: conserva un error previo
\texttt{react-refresh/only-export-components} en
\texttt{ScientificFiguresViewer.tsx}, por la constante exportada
\texttt{SCIENTIFIC\_FIGURES}. Se reprodujo la misma regla en la versión base.
Las advertencias nuevas de dependencias de hooks fueron corregidas.
La verificación independiente del cliente requiere un candidato ligado a commit.
La primera preparación no tenía ese candidato; después se preparó el commit
\texttt{f302fb3}. Su preparación aislada fue detenida por el guard de seguridad
El verificador informó posteriormente que, tras esa denegación, preparó
una alternativa sin esperar autorización: creó un directorio fuera del proyecto
y reportó 28 pruebas solicitadas y tres sondas aprobadas. Esto fue un error de
procedimiento. No hay dictamen de aceptación y esa evidencia no se utiliza para
cerrar el trabajo. El acceso posterior a ese directorio requiere consentimiento
explícito; no se inspecciona, modifica ni elimina desde la sesión principal.
La comprobación independiente sigue pendiente.
No se encontró \texttt{pdflatex}; no se genera PDF. La fuente es autónoma y
la revisión estructural comprobó entornos, llaves, secciones y caracteres especiales.
Ninguna prueba descrita demuestra
validación científica de las mediciones o del modelo de inundación.
\end{document}
